\documentclass[a4paper,12pt]{article}
\usepackage[T2A]{fontenc}
\usepackage[utf8]{inputenc}
\usepackage[russian]{babel}
\usepackage{amsmath,amssymb,mathtools}
\renewcommand{\familydefault}{\sfdefault}
\usepackage{icomma}
\usepackage[left=19mm,right=19mm,top=18mm,bottom=19mm]{geometry}
\usepackage{microtype}
\usepackage{tikz}
\usetikzlibrary{calc,arrows.meta,positioning,shapes.geometric}
\usepackage{tkz-euclide}
\usepackage{enumitem,tabularx,booktabs}
\usepackage{fancyhdr,parskip,setspace,needspace}
\usepackage{xcolor,hyperref}
\definecolor{accent}{HTML}{166D72}
\definecolor{darkaccent}{HTML}{15484F}
\definecolor{textgray}{HTML}{31383A}
\definecolor{softgray}{HTML}{667579}
\definecolor{amber}{HTML}{AC6A38}
\hypersetup{colorlinks=true,linkcolor=darkaccent,urlcolor=darkaccent}
\setlength{\parindent}{0pt}
\setlength{\parskip}{4pt}
\setstretch{1.18}
\setlist[itemize]{leftmargin=7mm,itemsep=2pt,topsep=3pt}
\setlist[enumerate]{leftmargin=9mm,itemsep=5pt,topsep=3pt}
\pagestyle{fancy}
\fancyhf{}
\fancyfoot[C]{\small\color{softgray}\thepage}
\renewcommand{\headrulewidth}{0pt}
\newcommand{\sect}[1]{\par\vspace{2mm}{\Large\bfseries\color{darkaccent}#1}\par\vspace{1mm}{\color{accent!50}\rule{\linewidth}{.55pt}}\par\vspace{1mm}}
\newcommand{\key}[1]{\textbf{\color{darkaccent}#1}}
\newcommand{\ck}[1]{\par\noindent$\square$\enspace #1\par}
\begin{document}
\thispagestyle{empty}
\begin{center}
\vspace*{33mm}
{\small\color{softgray}ГЕОМЕТРИЯ \textbullet{} ОГЭ \textbullet{} ЗАДАЧИ ПОВЫШЕННОЙ СЛОЖНОСТИ}\par
\vspace{14mm}
{\fontsize{30}{34}\selectfont\bfseries\color{darkaccent}Чек-лист}\par
\vspace{7mm}
{\LARGE\bfseries Отношения отрезков}\par
\vspace{3mm}
{\Large Теорема Менелая и вспомогательные параллельные прямые}\par
\vspace{19mm}
{\large София Кальней}\par
\vspace{4mm}
{\normalsize 16 октября 2026 года}\par
\vfill
{\small\color{textgray}Лёвин Артём Александрович эксклюзивно для Софии Кальней}\par
\vspace{17mm}
\end{center}
\begin{tikzpicture}[remember picture,overlay]
\draw[accent!50,line width=1.2pt]
 ([xshift=9mm,yshift=8mm]current page.south west)
 .. controls +(45mm,23mm) and +(-54mm,-6mm) ..
 ([xshift=-9mm,yshift=16mm]current page.south east);
\draw[amber!45,line width=.65pt]
 ([xshift=15mm,yshift=6mm]current page.south west)
 .. controls +(52mm,18mm) and +(-65mm,-8mm) ..
 ([xshift=-15mm,yshift=12mm]current page.south east);
\end{tikzpicture}
\clearpage

\sect{1. Главная идея: получить отношение из подобия}
\ck{Отличаю \textbf{часть отрезка} от \textbf{полного отрезка}: если $A,C,N$ лежат на одной прямой в порядке $A-C-N$, то $AN=AC+CN$.}
\ck{Нахожу сходственные стороны подобных треугольников и составляю отношения \emph{в одном порядке}.}
\ck{Понимаю, для чего проводится вспомогательная прямая, параллельная стороне треугольника.}

\key{Идея вспомогательной параллельной прямой.} Если задача содержит секущую и отношения отрезков, проведите через подходящую вершину прямую, параллельную одной из сторон. Пересечение с секущей часто создаёт \textbf{две связанные пары подобных треугольников}. Название «прямая Евклида», использовавшееся на уроке, условное; математический приём --- вспомогательное параллельное построение.

\sect{2. Теорема Менелая}
Пусть в треугольнике $ABC$ точки $K\in AB$, $J\in BC$, а точка $N$ находится \emph{на продолжении} $AC$ за $C$. Если точки $K,J,N$ лежат на одной прямой, то
\[
\boxed{\frac{AK}{KB}\cdot\frac{BJ}{JC}\cdot\frac{CN}{NA}=1.}
\]
\begin{center}
\begin{tikzpicture}[scale=1.08]
\tkzDefPoint(0,0){A}
\tkzDefPoint(.7,2){B}
\tkzDefPoint(2,0){C}
\tkzDefPoint(4,0){N}
\tkzDefPoint(.35,1){K}
\tkzDefPoint(1.566667,.666667){J}
\tkzDrawSegments[draw=textgray,line width=1pt](A,B B,C A,N)
\tkzDrawSegments[draw=accent,line width=1.35pt](K,N)
\tkzDrawPoints[size=2.5,fill=darkaccent](A,B,C,N,K,J)
\tkzLabelPoint[below left](A){$A$}
\tkzLabelPoint[above](B){$B$}
\tkzLabelPoint[below](C){$C$}
\tkzLabelPoint[below](N){$N$}
\tkzLabelPoint[left](K){$K$}
\tkzLabelPoint[right](J){$J$}
\end{tikzpicture}
\end{center}
\key{Как запомнить.} Идите по трём сторонам треугольника по кругу. На каждой стороне записывайте \emph{отношение двух последовательных отрезков}. Если одна точка внешняя, её отношение обязательно содержит \textbf{полную продлённую сторону}.

\key{Вторая конфигурация.} Если $K$ лежит на продолжении $AB$ за $B$, $M\in BC$, $N\in AC$ и $K,M,N$ коллинеарны, то
\[
\boxed{\frac{AN}{NC}\cdot\frac{CM}{MB}\cdot\frac{BK}{KA}=1,\qquad KA=AB+BK.}
\]
\key{Важно.} В этих формулах используются обычные положительные длины отрезков. Перемножать отношения следует в согласованном порядке, сохраняя полные внешние отрезки $NA$ или $KA$.
\clearpage

\sect{3. Задача с занятия: середина и продолжение стороны}
\textbf{Дано:} $K$ --- середина $AB$; на продолжении $AC$ за $C$ взята точка $N$ так, что $CN=AC$. Прямая $KN$ пересекает $BC$ в точке $J$. Найти $BJ:JC$.

\begin{center}
\begin{tikzpicture}[scale=1.34]
\tkzDefPoint(0,0){A}
\tkzDefPoint(.7,2){B}
\tkzDefPoint(2,0){C}
\tkzDefPoint(4,0){N}
\tkzDefPoint(.35,1){K}
\tkzDefPoint(1.566667,.666667){J}
\tkzDefPoint(-3.3,2){F}
\tkzDrawSegments[draw=textgray,line width=.95pt](A,B B,C A,N)
\tkzDrawSegments[draw=accent,line width=1.15pt](F,B F,N)
\tkzDrawPoints[size=2.2,fill=darkaccent](A,B,C,N,K,J,F)
\tkzLabelPoint[below](A){$A$}
\tkzLabelPoint[above](B){$B$}
\tkzLabelPoint[below](C){$C$}
\tkzLabelPoint[below](N){$N$}
\tkzLabelPoint[left](K){$K$}
\tkzLabelPoint[right](J){$J$}
\tkzLabelPoint[above left](F){$F$}
\node[font=\scriptsize,text=accent] at (-1.2,2.24) {$FB\parallel AN$};
\end{tikzpicture}
\end{center}

\key{Способ 1. Вспомогательная параллельная.} Проведём через $B$ прямую $BF\parallel AC$, где $F\in KN$.
\[
\triangle FKB\sim\triangle NKA
\quad\Longrightarrow\quad
\frac{FB}{AN}=\frac{KB}{KA}=1.
\]
Так как $AN=AC+CN=2CN$, то $FB=2CN$. Кроме того,
\[
\triangle FBJ\sim\triangle NCJ
\quad\Longrightarrow\quad
\frac{BJ}{JC}=\frac{FB}{CN}=2.
\]
\key{Способ 2. Теорема Менелая.}
\[
\underbrace{\frac{AK}{KB}}_{1}\cdot\frac{BJ}{JC}\cdot
\underbrace{\frac{CN}{NA}}_{1/2}=1
\quad\Longrightarrow\quad\boxed{BJ:JC=2:1.}
\]
\key{Смысл сравнения.} Подобие показывает, \emph{почему} возникает искомое отношение; теорема Менелая позволяет получить его одной цепочкой вычислений.
\clearpage

\sect{4. Задача с занятия: внешняя точка на другой стороне}
\textbf{Дано:} $K$ лежит на продолжении $AB$ за $B$; $M\in BC$; прямая $KM$ пересекает $AC$ в точке $N$; $BK:AB=1:5$, $BM:MC=4:5$. Найти $CN:AN$.

\begin{center}
\begin{tikzpicture}[scale=1.35]
\tkzDefPoint(0,0){A}
\tkzDefPoint(4,0){B}
\tkzDefPoint(1.2,4){C}
\tkzDefPoint(4.8,0){K}
\tkzDefPoint(2.755556,1.777778){M}
\tkzDefPoint(.993103,3.310345){N}
\tkzDefPoint(.2,4){J}
\tkzDrawSegments[draw=textgray,line width=.95pt](A,K B,C A,C)
\tkzDrawSegments[draw=accent,line width=1.15pt](K,J C,J)
\tkzDrawPoints[size=2.25,fill=darkaccent](A,B,C,K,M,N,J)
\tkzLabelPoint[below](A){$A$}
\tkzLabelPoint[below](B){$B$}
\tkzLabelPoint[above](C){$C$}
\tkzLabelPoint[below](K){$K$}
\tkzLabelPoint[right](M){$M$}
\tkzLabelPoint[left](N){$N$}
\tkzLabelPoint[above](J){$J$}
\node[font=\scriptsize,text=accent] at (.7,4.28) {$CJ\parallel AK$};
\end{tikzpicture}
\end{center}

\key{Подготовка.} $BK:AB=1:5$. Поскольку $AK=AB+BK$, получаем $BK:AK=1:6$. В отношениях $BM:MC=4:5$ и $BK:AB=1:5$ единицу масштаба \emph{нельзя автоматически считать одной и той же}.

\key{Способ 1. Две пары подобных треугольников.} Проведём $CJ\parallel AB$, $J\in KM$.
\[
\triangle KBM\sim\triangle JCM
\ \Longrightarrow\ \frac{CJ}{BK}=\frac{CM}{BM}=\frac54.
\]
\[
\triangle KAN\sim\triangle JCN
\ \Longrightarrow\ \frac{CN}{AN}=\frac{CJ}{KA}
=\frac54\cdot\frac16=\boxed{\frac5{24}}.
\]
\key{Способ 2. Теорема Менелая.}
\[
\frac{AN}{NC}\cdot\underbrace{\frac{CM}{MB}}_{5/4}
\cdot\underbrace{\frac{BK}{KA}}_{1/6}=1
\Longrightarrow \frac{AN}{NC}=\frac{24}{5}.
\]
Следовательно, \(\boxed{CN:AN=5:24}\).
\clearpage

\sect{5. Что нужно проверять в каждом решении}
\ck{На чертеже обозначены \textbf{все} точки пересечения и порядок точек на прямых.}
\ck{Я различаю $CN:AN$ и $CN:AC$, а также $BK:KA$ и $BK:BA$.}
\ck{В теореме Менелая использованы три отношения, произведение которых равно $1$.}
\ck{В подобии сопоставлены \emph{сходственные} стороны.}
\ck{При записи $AB=5x$ и $BM=4y$ использованы разные масштабы $x$ и $y$, пока не доказана связь между ними.}
\ck{В конце дан ответ в том порядке, который указан в вопросе.}

\key{Типичные ошибки.}
\begin{itemize}
\item Подставляют $AC$ вместо $AN=AC+CN$ или $AB$ вместо $AK=AB+BK$.
\item Переворачивают только одно отношение, нарушая согласованный обход треугольника.
\item Принимают пропорции разных сторон за равенство реальных длин.
\item Путают порядок отношений: ответ $AN:CN$ заменяют на $CN:AN$ без обращения дроби.
\end{itemize}

\sect{6. Как выбрать метод}
\textbf{Теорема Менелая} особенно удобна, когда секущая пересекает три стороны треугольника (в том числе их продолжения) и надо вычислить неизвестное отношение.

\textbf{Параллельное вспомогательное построение} полезно, когда требуется понять происхождение теоремы или найти связанные пары подобных треугольников. Оно также помогает увидеть, какие стороны являются сходственными.

\key{Самопроверка.} Полученные отношения должны быть положительными; если точка лежит \emph{внутри} стороны, соответствующая доля этой стороны строго меньше $1$. Например, для $N\in AC$ всегда $CN/AN>0$ без универсальной верхней границы, а для $N$ на продолжении $AC$ за $C$ всегда $CN/AN<1$.

\clearpage
\begin{center}
{\Large\bfseries\color{darkaccent}Алгоритм: найти отношение с помощью теоремы Менелая}\par
\vspace{4mm}
{\small\color{softgray}Сначала распознать конфигурацию, затем составить три отношения}
\end{center}
\vspace{2mm}
\begin{center}
\begin{tikzpicture}[
  >=Stealth, node distance=9mm,
  start/.style={ellipse,draw=darkaccent,line width=.85pt,minimum height=10mm,align=center,text width=7.0cm,fill=accent!4},
  flowstep/.style={rectangle,rounded corners=1.3mm,draw=accent!85,line width=.8pt,minimum height=12mm,align=center,text width=10.6cm,inner sep=3mm},
  choose/.style={diamond,draw=amber,line width=.9pt,aspect=2.65,align=center,inner sep=2mm,text width=5cm},
  branch/.style={rectangle,rounded corners=1mm,draw=accent!75,line width=.8pt,minimum height=22mm,align=center,text width=5.2cm,inner sep=2mm},
  arrow/.style={->,draw=darkaccent,line width=1pt}
]
\node[start] (s) {Треугольник и прямая, пересекающая три его стороны или продолжения};
\node[flowstep,below=of s] (a) {Найдите на рисунке \textbf{внешнюю точку} и определите полные стороны};
\node[choose,below=11mm of a] (d) {Внешняя точка на продолжении $AC$?};
\node[branch,below=17mm of d,xshift=-36mm] (l) {$K\in AB$, $J\in BC$, $N$ за $C$:\\[1mm] $\displaystyle\frac{AK}{KB}\cdot\frac{BJ}{JC}\cdot\frac{CN}{NA}=1$};
\node[branch,below=17mm of d,xshift=36mm] (r) {$K$ за $B$, $M\in BC$, $N\in AC$:\\[1mm] $\displaystyle\frac{AN}{NC}\cdot\frac{CM}{MB}\cdot\frac{BK}{KA}=1$};
\coordinate (mid) at ($(l.south)!0.5!(r.south)$);
\node[flowstep,below=16mm of mid] (b) {Подставьте известные отношения, найдите неизвестное};
\node[flowstep,below=of b] (c) {Сверьте порядок отношения с вопросом; проверьте полные отрезки};
\node[start,below=of c] (e) {Ответ};
\draw[arrow] (s)--(a);
\draw[arrow] (a)--(d);
\draw[arrow] (d.south) -- ++(0,-5mm) -| node[pos=.87,above,font=\small] {да} (l.north);
\draw[arrow] (d.south) -- ++(0,-5mm) -| node[pos=.87,above,font=\small] {нет} (r.north);
\coordinate (join) at ($(b.north)+(0,5mm)$);
\draw[draw=darkaccent,line width=1pt] (l.south) |- (join);
\draw[draw=darkaccent,line width=1pt] (r.south) |- (join);
\draw[arrow] (join)--(b.north);
\draw[arrow] (b)--(c);
\draw[arrow] (c)--(e);
\end{tikzpicture}
\end{center}
\vfill
{\small\color{softgray}Для иных расположений точек заново запишите произведение отношений по согласованному обходу треугольника.}
\clearpage

\sect{7. Самостоятельная тренировка --- Путь А}
\textbf{10 заданий.} Чертёж к геометрической задаче выполняйте самостоятельно. Если используется теорема Менелая, записывайте все три множителя.
\begin{enumerate}[label=\textbf{\arabic*.},leftmargin=9mm,itemsep=5pt]
\item Точки $A,C,N$ лежат на одной прямой в порядке $A-C-N$: $AC=6$, $CN=9$. На продолжении $AB$ за $B$ взята $K$: $AB=12$, $BK=4$. Найдите $AN$ и $AK$.
\item В $\triangle ABC$ точка $K\in AB$, точка $N$ --- на продолжении $AC$ за $C$, а $J=KN\cap BC$. Известно $AK:KB=2:3$, $CN:AC=1:2$. Найдите $BJ:JC$.
\item В той же конфигурации $K$ --- середина $AB$, а $CN=AC$. Найдите $BJ:JC$.
\item В той же конфигурации $K$ --- середина $AB$, $CN:AC=1:2$ и $BC=24$. Найдите $BJ$ и $JC$.
\item В той же конфигурации $AK:KB=3:2$, $BJ:JC=5:1$. Найдите $CN:AC$.
\item В той же конфигурации $BJ:JC=6:1$, $CN:AC=1:2$. Найдите $AK:KB$.
\item В $\triangle ABC$ точка $K$ находится на продолжении $AB$ за $B$, $M\in BC$, $N=KM\cap AC$. Если $BK:AB=1:3$, $BM:MC=2:5$, найдите $CN:AN$.
\item Для той же конфигурации $AB=18$, $BK=6$, $BM=6$, $MC=9$. Найдите $CN:AN$.
\item В первой конфигурации $K$ --- середина $AB$; $AC:CN=3:2$. Через $B$ проведена прямая $BF\parallel AC$, где $F\in KN$. Найдите $BF:CN$ и $BJ:JC$ через подобие.
\item В конфигурации задания 7 дано $BK:AB=2:7$, $BM:MC=3:8$. Найдите $CN:AN$ \textbf{двумя способами}: через Менелая и с помощью вспомогательной прямой через $C$, параллельной $AB$.
\end{enumerate}
\clearpage

\sect{8. Ответы и ориентиры для самопроверки}
\begin{enumerate}[label=\textbf{\arabic*.},leftmargin=9mm,itemsep=7pt]
\item $AN=15$, $AK=16$. Внешняя сторона равна сумме своих частей.
\item $BJ:JC=9:2$, поскольку $\frac{AK}{KB}=\frac23$, $\frac{CN}{NA}=\frac13$.
\item $BJ:JC=2:1$.
\item $BJ=18$, $JC=6$: отношение $BJ:JC=3:1$, а $BJ+JC=24$.
\item $CN:AC=2:13$. Сначала $CN:AN=2:15$; затем $AN=AC+CN$.
\item $AK:KB=1:2$: $\frac{AK}{KB}\cdot 6\cdot\frac13=1$.
\item $CN:AN=5:8$: $\frac{BK}{KA}=\frac14$, $\frac{CM}{MB}=\frac52$, откуда $\frac{AN}{NC}=\frac85$.
\item $CN:AN=3:8$: $\frac{BK}{KA}=\frac14$ и $\frac{CM}{MB}=\frac32$.
\item $BF:CN=5:2$; $BJ:JC=5:2$. Здесь $AN:CN=5:2$, а $BK=AK$.
\item $CN:AN=16:27$. Менелай: $\frac{AN}{NC}\cdot\frac83\cdot\frac29=1$. Через подобие: $\frac{CJ}{BK}=\frac83$ и $\frac{CN}{AN}=\frac{CJ}{KA}$.
\end{enumerate}
\vfill
{\small\color{softgray}Для заданий 2--10 полезно сверить каждое вычисление с чертежом: где часть, где полная сторона и каков порядок точек.}
\end{document}
